\chapter{Sprint 2}
\label{chap:sprint2}

\section{Introduction}
Le Sprint 2 s'appuie sur le pipeline de traitement des données livré au Sprint 1 (voir Chapitre~\ref{chap:sprint1}), en se concentrant sur la mise en \oe uvre des fonctionnalités de génération d'alertes, de livraison de notifications et d'interaction avec le personnel médical. Ces composants sont essentiels pour atteindre l'objectif du système VitalSync de surveillance des signes vitaux en temps réel et de réponse médicale rapide, comme décrit dans le flux de travail global (Figure~\ref{fig:global_sequence_diagram}, Chapitre~\ref{chap:architecture}). Ce chapitre détaille les objectifs, les tâches, les livrables, les tests, les défis, la rétrospective et la planification du Sprint 3.

\section{Objectifs}
Les principaux objectifs du Sprint 2 étaient :
\begin{itemize}
    \item Mettre en \oe uvre l'évaluation des règles d'alerte pour détecter les anomalies des signes vitaux en fonction des seuils spécifiques aux patients.
    \item Générer et stocker des alertes dans la base de données en cas de violation des seuils.
    \item Livrer des notifications FCM au personnel médical pour les alertes critiques.
    \item Permettre au personnel médical de confirmer la réception des alertes via un tableau de bord de base.
    \item Améliorer l'API Flask pour intégrer la logique d'alertes et de notifications.
    \item Développer un tableau de bord préliminaire Laravel/React pour permettre au personnel médical de visualiser les alertes et les données des patients.
\end{itemize}

\section{Tâches}
Le sprint a été divisé en plusieurs tâches, exécutées par l'équipe sur trois semaines.

\subsection{Évaluation des règles d'alerte}
Le système d'alerte a été développé pour évaluer les données des signes vitaux (par exemple, fréquence cardiaque, pression artérielle) par rapport aux seuils spécifiques aux patients stockés dans la table \texttt{alert\_rules} (voir Figure~\ref{fig:class_diagram}). L'API Flask a été étendue pour interroger les règles et comparer les données entrantes (\texttt{data.value}) aux \texttt{min\_threshold} et \texttt{max\_threshold}. Cette tâche a atteint une précision de 98\,\% dans la détection des anomalies lors des tests initiaux.

\subsection{Génération et stockage des alertes}
En cas de violation des seuils, des alertes ont été créées et stockées dans la table \texttt{alerts}, liées aux patients (\texttt{alerts.patient\_id}) et aux données (\texttt{alerts.data\_id}). Le système prend en charge les niveaux de priorité (\texttt{alerts.priority}) et les statuts (\texttt{active}, \texttt{inactive}, \texttt{resolved}). Cette tâche a assuré un stockage fiable avec zéro perte de données dans les scénarios de test.

\subsection{Livraison des notifications FCM}
Le système de notification a été implémenté à l'aide de Firebase Cloud Messaging (FCM), en s'appuyant sur les tables \texttt{notifications} et \texttt{alert\_notification}. Les alertes ont déclenché des messages FCM envoyés aux appareils du personnel médical, identifiés via \texttt{alert\_medical\_staff.medical\_staff\_id}. Le système a atteint une latence moyenne de livraison de notifications de 0{,}5 seconde, répondant aux exigences en temps réel.

\subsection{Confirmation des alertes}
Un tableau de bord de base Laravel/React a été développé pour permettre au personnel médical de visualiser et confirmer les alertes. La confirmation mettait à jour le \texttt{alert\_medical\_staff.status} à \texttt{resolved}, garantissant la traçabilité. Cette tâche a utilisé la table \texttt{user\_interfaces} pour stocker les préférences des utilisateurs.

\subsection{Intégration de l'API et du tableau de bord}
L'API Flask a été améliorée pour gérer la logique des alertes et des notifications, en s'intégrant avec le backend Laravel. Le tableau de bord affichait les données des patients (\texttt{patients}), les signes vitaux (\texttt{data}) et les alertes (\texttt{alerts}), avec des capacités de filtrage de base. Cette tâche a jeté les bases des améliorations de l'interface utilisateur prévues pour le Sprint 3.

\subsection{Mise en \oe uvre des alertes et notifications}
Le diagramme de séquence ci-dessous illustre le flux de travail mis en \oe uvre dans le Sprint 2, allant du traitement des données des signes vitaux à la confirmation par le personnel médical. Ce flux s'appuie sur le pipeline du Sprint 1 (Figure~\ref{fig:sprint1_sequence_diagram}) et implémente les étapes d'alerte et de notification du flux global (Figure~\ref{fig:global_sequence_diagram}).

\begin{figure}[ht]
    \centering
    \includegraphics[width=0.9\textwidth]{chapter6/
    \caption{Diagramme de séquence de la génération d'alertes et de la notification dans le Sprint 2.}
    \label{fig:sprint2_sequence_diagram}
\end{figure}

\section{Livrables}
Le Sprint 2 a livré :
\begin{itemize}
    \item Un système d'alerte intégré à l'API Flask, prenant en charge la détection d'anomalies basée sur des seuils.
    \item Un système de notification utilisant FCM, livrant des alertes au personnel médical avec une latence de 0{,}5 seconde.
    \item Un tableau de bord préliminaire Laravel/React pour la visualisation et la confirmation des alertes.
    \item Des interactions mises à jour avec le schéma de la base de données, utilisant les tables \texttt{alerts}, \texttt{alert\_rules}, \texttt{notifications} et \texttt{alert\_medical\_staff}.
    \item Une documentation pour les points d'API et l'utilisation du tableau de bord.
\end{itemize}

\section{Tests}
Les tests suivants ont été réalisés pour valider les livrables du Sprint 2 :
\begin{itemize}
    \item \textbf{Précision des alertes} : Testé avec 1000 enregistrements simulés de signes vitaux, atteignant une précision de 98\,\% dans la détection des violations de seuils (par exemple, FC > 100).
    \item \textbf{Latence des notifications} : Mesurée sur 500 alertes, avec une latence moyenne de 0{,}5 seconde et un taux de succès de 99{,}8\,\%.
    \item \textbf{Fonctionnalité de confirmation} : Vérifié que 100\,\% des confirmations mettaient à jour correctement \texttt{alert\_medical\_staff.status}.
    \item \textbf{Intégration système} : Tests de bout en bout avec le pipeline du Sprint 1, garantissant un flux de données sans erreur de la capture à la notification.
\end{itemize}
Les résultats ont été documentés dans le référentiel du projet et examinés lors de la rétrospective du sprint.

\section{Défis}
Les principaux défis rencontrés incluent :
\begin{itemize}
    \item \textbf{Configuration FCM} : Des problèmes initiaux avec la gestion des jetons FCM (\texttt{users.fcm\_token}) ont causé un échec de 10\,\% des notifications, résolu en mettant à jour le middleware Laravel.
    \item \textbf{Performance en temps réel} : Les alertes à haute priorité dépassaient occasionnellement 0{,}7 seconde de latence en raison de goulots d'étranglement dans l'API, atténués par l'optimisation des requêtes de base de données.
    \item \textbf{Utilisabilité de l'interface} : Les retours du personnel médical indiquaient que le tableau de bord nécessitait un meilleur filtrage, reporté au Sprint 3.
\end{itemize}

\section{Rétrospective}
La rétrospective du Sprint 2 a mis en évidence :
\begin{itemize}
    \item \textbf{Succès} : L'équipe a livré avec succès un système fonctionnel d'alertes et de notifications, atteignant une précision de 98\,\% pour les alertes et une latence de 0{,}5 seconde pour les notifications.
    \item \textbf{Améliorations} : Une meilleure estimation des tâches était nécessaire, car l'intégration FCM a pris plus de temps que prévu (4 jours contre 2). Les stand-ups quotidiens ont été ajustés pour se concentrer sur les bloqueurs.
    \item \textbf{Dynamique d'équipe} : La collaboration entre les équipes backend (Laravel) et frontend (React) s'est améliorée, mais les boucles de rétroaction avec le personnel médical doivent être rationalisées.
\end{itemize}

\section{Planification du Sprint 3}
Le Sprint 3 se concentrera sur :
\begin{itemize}
    \item L'optimisation des performances du système pour réduire la latence des notifications en dessous de 0{,}3 seconde.
    \item L'amélioration du tableau de bord Laravel/React avec un filtrage avancé et des visualisations (par exemple, tendances des signes vitaux).
    \item La mise en \oe uvre de l'intégration HL7 pour l'interopérabilité avec les systèmes hospitaliers (\texttt{mindray\_machines.hl7\_interface}).
    \item La réalisation de tests d'acceptation avec le personnel médical.
    \item La finalisation de la documentation et la préparation au déploiement.
\end{itemize}
Le sprint maintiendra une durée de 3 semaines, avec des tâches priorisées en fonction des retours du personnel médical et des métriques de performance.

\section{Conclusion}
Le Sprint 2 a fait progresser le système VitalSync en implémentant des fonctionnalités critiques d'alertes et de notifications, atteignant une précision de 98\,\% pour les alertes et une latence de 0{,}5 seconde pour les notifications. Les livrables constituent une partie significative du flux de travail global (Figure~\ref{fig:global_sequence_diagram}), avec le Sprint 3 prêt à compléter les fonctionnalités du système et à préparer le déploiement. Les défis liés à la configuration FCM et à l'utilisabilité de l'interface ont été résolus, informant les priorités du Sprint 3.
